% aspect ratio: 16:9
\documentclass[unicode, 9pt, aspectratio=169]{beamer}
\usepackage{mybeamer}

%%%%%%%%%%%%%%%%%%%%%%%%%
% document
%%%%%%%%%%%%%%%%%%%%%%%%%
% title
\title{鉄鋼スラグを活用したカルシア改質土の\\実用化に向けた土質特性の評価と品質管理に関する研究}
\author{博士論文発表}
\institute{京都大学大学院 工学研究科}
\date{\today}

% main
\begin{document}
\newcommand{\NoTranslate}[1]{#1}
\pdfstringdefDisableCommands{\let\translate\NoTranslate}

\maketitle

%%%%%%%%%%%%%%%%%%%%%%%%%
% Introduction
%%%%%%%%%%%%%%%%%%%%%%%%%
\section{研究背景と目的}

\begin{frame}
  \frametitle{研究の背景}
  \begin{columns}
    \begin{column}{0.55\linewidth}
      \begin{block}{社会的課題}
        \begin{itemize}
          \item 軟弱地盤の改良需要が増大
          \item 建設残土の有効活用が課題
          \item 環境負荷の低減が求められる
        \end{itemize}
      \end{block}
      
      \begin{alertblock}{鉄鋼スラグの可能性}
        \begin{itemize}
          \item 年間約4000万トンの鉄鋼スラグが発生
          \item カルシウム成分が豊富
          \item 土の改質材料として期待
        \end{itemize}
      \end{alertblock}
    \end{column}
    
    \begin{column}{0.40\linewidth}
      \begin{exampleblock}{本研究の意義}
        \centering
        \Large
        \emph{持続可能な}\\
        \emph{地盤改良技術}\\
        \vspace{0.5em}
        の確立
      \end{exampleblock}
      
      \vspace{1em}
      \begin{minipage}{\linewidth}
        \centering
        \small
        鉄鋼スラグ $\to$ 土壌改質\\
        $\downarrow$\\
        環境配慮型建設材料
      \end{minipage}
    \end{column}
  \end{columns}
\end{frame}

\begin{frame}
  \frametitle{研究目的}
  \begin{center}
    \begin{minipage}{0.85\linewidth}
      \begin{block}{主要な研究目的}
        \begin{enumerate}
          \item 鉄鋼スラグを活用したカルシア改質土の\emph{土質特性}の解明
          \item 改質土の\emph{強度発現メカニズム}の解明
          \item 実用化に向けた\emph{品質管理手法}の確立
          \item 環境安全性の評価
        \end{enumerate}
      \end{block}
    \end{minipage}
  \end{center}
  
  \vspace{1em}
  
  \begin{columns}
    \begin{column}{0.48\linewidth}
      \begin{roundbox}[blue!20]{0.9\linewidth}
        \centering
        \small
        \textbf{学術的貢献}\\
        改質メカニズムの\\
        科学的解明
      \end{roundbox}
    \end{column}
    \begin{column}{0.48\linewidth}
      \begin{roundbox}[green!20]{0.9\linewidth}
        \centering
        \small
        \textbf{実用的貢献}\\
        品質管理基準の\\
        策定
      \end{roundbox}
    \end{column}
  \end{columns}
\end{frame}

%%%%%%%%%%%%%%%%%%%%%%%%%
% Materials and Methods
%%%%%%%%%%%%%%%%%%%%%%%%%
\section{実験方法と材料}

\begin{frame}
  \frametitle{使用材料}
  \begin{columns}
    \begin{column}{0.48\linewidth}
      \begin{block}{土質材料}
        \begin{description}[含水比]
          \item[種類] 沖積粘性土
          \item[液性限界] 52.3\%
          \item[塑性限界] 28.5\%
          \item[含水比] 45.2\%
        \end{description}
      \end{block}
      
      \begin{block}{鉄鋼スラグ}
        \begin{description}[粒度]
          \item[種類] 製鋼スラグ
          \item[主成分] CaO, SiO$_2$, Fe$_2$O$_3$
          \item[粒度] 75$\mu$m以下
        \end{description}
      \end{block}
    \end{column}
    
    \begin{column}{0.48\linewidth}
      \begin{alertblock}{配合条件}
        \begin{itemize}
          \item スラグ添加率：0, 5, 10, 15\%
          \item 養生温度：20$^\circ$C
          \item 養生期間：0〜91日
          \item 含水比調整：最適含水比
        \end{itemize}
      \end{alertblock}
      
      \vspace{1em}
      
      \begin{exampleblock}{}
        \centering
        \small
        \emph{体系的な}\\
        \emph{実験計画}
      \end{exampleblock}
    \end{column}
  \end{columns}
\end{frame}

\begin{frame}
  \frametitle{試験項目}
  \begin{center}
    \begin{minipage}{0.85\linewidth}
      \begin{block}{実施した試験}
        \begin{columns}
          \begin{column}{0.48\linewidth}
            \textbf{力学特性試験}
            \begin{itemize}
              \item 一軸圧縮試験
              \item CBR試験
              \item 三軸圧縮試験
            \end{itemize}
            
            \vspace{0.5em}
            \textbf{物理特性試験}
            \begin{itemize}
              \item pH測定
              \item 密度試験
              \item 透水試験
            \end{itemize}
          \end{column}
          
          \begin{column}{0.48\linewidth}
            \textbf{化学分析}
            \begin{itemize}
              \item XRD分析
              \item SEM観察
              \item 溶出試験
            \end{itemize}
            
            \vspace{0.5em}
            \textbf{長期特性}
            \begin{itemize}
              \item 耐久性試験
              \item 環境安全性評価
            \end{itemize}
          \end{column}
        \end{columns}
      \end{block}
    \end{minipage}
  \end{center}
\end{frame}

%%%%%%%%%%%%%%%%%%%%%%%%%
% Results
%%%%%%%%%%%%%%%%%%%%%%%%%
\section{実験結果}

\begin{frame}
  \frametitle{一軸圧縮強度の経時変化}
  \begin{center}
    \begin{tikzpicture}
      \begin{axis}[
        width=0.85\textwidth,
        height=0.6\textheight,
        xlabel={養生日数 (日)},
        ylabel={一軸圧縮強度 (kN/m$^2$)},
        xmin=0, xmax=100,
        ymin=0, ymax=750,
        legend pos=north west,
        grid=major,
        grid style={dashed,gray!30},
        legend style={font=\small},
        tick label style={font=\small},
        label style={font=\normalsize}
      ]
      
      \addplot[color=blue,mark=*,thick] table[x=curing_days,y=unconfined_strength_control,col sep=comma] {data/strength_data.csv};
      \addlegendentry{無添加}
      
      \addplot[color=green,mark=square*,thick] table[x=curing_days,y=unconfined_strength_slag5,col sep=comma] {data/strength_data.csv};
      \addlegendentry{スラグ5\%}
      
      \addplot[color=orange,mark=triangle*,thick] table[x=curing_days,y=unconfined_strength_slag10,col sep=comma] {data/strength_data.csv};
      \addlegendentry{スラグ10\%}
      
      \addplot[color=red,mark=diamond*,thick] table[x=curing_days,y=unconfined_strength_slag15,col sep=comma] {data/strength_data.csv};
      \addlegendentry{スラグ15\%}
      
      \end{axis}
    \end{tikzpicture}
  \end{center}
  
  \begin{minipage}{0.6\linewidth}
    \begin{roundbox}[yellow!30]{\linewidth}
      \centering
      \small
      スラグ添加により強度が\emph{最大2.2倍}に向上
    \end{roundbox}
  \end{minipage}
\end{frame}

\begin{frame}
  \frametitle{pH値の変化}
  \begin{columns}
    \begin{column}{0.55\linewidth}
      \begin{tikzpicture}
        \begin{axis}[
          width=\textwidth,
          height=0.65\textheight,
          xlabel={スラグ添加率 (\%)},
          ylabel={pH},
          xmin=-1, xmax=16,
          ymin=7, ymax=14,
          legend pos=north west,
          grid=major,
          grid style={dashed,gray!30},
          legend style={font=\scriptsize},
          tick label style={font=\scriptsize},
          label style={font=\small}
        ]
        
        \addplot[color=blue,mark=*,thick] table[x=slag_content,y=ph_initial,col sep=comma] {data/ph_data.csv};
        \addlegendentry{初期}
        
        \addplot[color=green,mark=square*,thick] table[x=slag_content,y=ph_7days,col sep=comma] {data/ph_data.csv};
        \addlegendentry{7日}
        
        \addplot[color=orange,mark=triangle*,thick] table[x=slag_content,y=ph_28days,col sep=comma] {data/ph_data.csv};
        \addlegendentry{28日}
        
        \addplot[color=red,mark=diamond*,thick] table[x=slag_content,y=ph_91days,col sep=comma] {data/ph_data.csv};
        \addlegendentry{91日}
        
        \end{axis}
      \end{tikzpicture}
    \end{column}
    
    \begin{column}{0.42\linewidth}
      \begin{alertblock}{重要な知見}
        \begin{itemize}
          \item スラグ添加でpHが上昇
          \item 高アルカリ環境の形成
          \item 養生日数とともにpH増加
        \end{itemize}
      \end{alertblock}
      
      \vspace{1em}
      
      \begin{exampleblock}{化学反応}
        \centering
        \small
        CaO + H$_2$O $\to$ Ca(OH)$_2$\\
        $\downarrow$\\
        強度発現の鍵
      \end{exampleblock}
    \end{column}
  \end{columns}
\end{frame}

\begin{frame}
  \frametitle{CBR値の改善効果}
  \begin{columns}
    \begin{column}{0.50\linewidth}
      \begin{tikzpicture}
        \begin{axis}[
          width=\textwidth,
          height=0.65\textheight,
          xlabel={スラグ添加率 (\%)},
          ylabel={CBR (\%)},
          xmin=-1, xmax=21,
          ymin=0, ymax=65,
          grid=major,
          grid style={dashed,gray!30},
          tick label style={font=\scriptsize},
          label style={font=\small},
          ybar,
          bar width=15pt,
          fill=blue!60,
          draw=blue!80,
        ]
        
        \addplot table[x=slag_content,y=cbr_value,col sep=comma] {data/cbr_data.csv};
        
        \end{axis}
      \end{tikzpicture}
    \end{column}
    
    \begin{column}{0.47\linewidth}
      \begin{block}{CBR試験結果}
        \begin{itemize}
          \item 無添加：8\%
          \item スラグ15\%：48\%
          \item \emph{6倍}の改善効果
        \end{itemize}
      \end{block}
      
      \vspace{1em}
      
      \begin{alertblock}{実用性評価}
        \centering
        \small
        路盤材として\\
        \emph{十分な支持力}\\
        を確保
      \end{alertblock}
      
      \vspace{1em}
      
      \begin{roundbox}[green!20]{\linewidth}
        \centering
        \scriptsize
        基準値（30\%）を\\
        スラグ10\%で達成
      \end{roundbox}
    \end{column}
  \end{columns}
\end{frame}

%%%%%%%%%%%%%%%%%%%%%%%%%
% Mechanism
%%%%%%%%%%%%%%%%%%%%%%%%%
\section{強度発現メカニズム}

\begin{frame}
  \frametitle{カルシア改質のメカニズム}
  \begin{center}
    \begin{minipage}{0.85\linewidth}
      \begin{block}{反応プロセス}
        \begin{enumerate}
          \item \textbf{初期段階}：スラグからのCaO溶出
          \item \textbf{水和反応}：Ca(OH)$_2$の生成
          \item \textbf{土粒子への作用}：カルシウムイオンの吸着
          \item \textbf{結晶生成}：C-S-H, C-A-H等の水和物形成
          \item \textbf{固化}：粒子間の結合強化
        \end{enumerate}
      \end{block}
    \end{minipage}
  \end{center}
  
  \vspace{0.5em}
  
  \begin{columns}
    \begin{column}{0.31\linewidth}
      \begin{roundbox}[blue!15]{0.95\linewidth}
        \centering
        \scriptsize
        \textbf{短期効果}\\
        イオン交換\\
        凝集作用
      \end{roundbox}
    \end{column}
    \begin{column}{0.31\linewidth}
      \begin{roundbox}[green!15]{0.95\linewidth}
        \centering
        \scriptsize
        \textbf{中期効果}\\
        ポゾラン反応\\
        水和物生成
      \end{roundbox}
    \end{column}
    \begin{column}{0.31\linewidth}
      \begin{roundbox}[red!15]{0.95\linewidth}
        \centering
        \scriptsize
        \textbf{長期効果}\\
        結晶成長\\
        緻密化
      \end{roundbox}
    \end{column}
  \end{columns}
\end{frame}

\begin{frame}
  \frametitle{微視的構造の変化}
  \begin{columns}
    \begin{column}{0.48\linewidth}
      \begin{alertblock}{SEM観察結果}
        \begin{itemize}
          \item 無添加土：粒子間に大きな空隙
          \item 改質土：水和物による充填
          \item 針状・板状結晶の形成
          \item 粒子間結合の強化
        \end{itemize}
      \end{alertblock}
      
      \vspace{1em}
      
      \begin{block}{XRD分析}
        検出された鉱物相
        \begin{itemize}
          \item C-S-H gel
          \item Ettringite
          \item Calcite
        \end{itemize}
      \end{block}
    \end{column}
    
    \begin{column}{0.48\linewidth}
      \begin{exampleblock}{構造変化の特徴}
        \centering
        \vspace{0.5em}
        \tikz{
          \node[draw, rectangle, fill=gray!30, minimum width=3cm, minimum height=1cm] (soil) {未改質土};
          \node[below=0.5cm of soil] (arrow) {$\downarrow$ スラグ添加};
          \node[draw, rectangle, fill=blue!20, minimum width=3cm, minimum height=1cm, below=0.5cm of arrow] (modified) {改質土};
        }
        \vspace{0.5em}
        
        \begin{itemize}
          \item 空隙率低減
          \item 密度増加
          \item 結合強度向上
        \end{itemize}
      \end{exampleblock}
    \end{column}
  \end{columns}
  
  \vspace{0.5em}
  \centering
  \begin{minipage}{0.6\linewidth}
    \begin{roundbox}[yellow!30]{\linewidth}
      \centering
      \small
      微視的構造変化が\emph{巨視的強度}に寄与
    \end{roundbox}
  \end{minipage}
\end{frame}

%%%%%%%%%%%%%%%%%%%%%%%%%
% Quality Control
%%%%%%%%%%%%%%%%%%%%%%%%%
\section{品質管理手法}

\begin{frame}
  \frametitle{品質管理基準の提案}
  \begin{center}
    \begin{minipage}{0.85\linewidth}
      \begin{block}{管理項目と基準値}
        \begin{center}
          \small
          \begin{tabular}{lcc}
            \toprule
            \textbf{管理項目} & \textbf{基準値} & \textbf{試験方法} \\
            \midrule
            スラグ添加率 & 10〜15\% & 重量比 \\
            含水比 & 最適含水比$\pm$2\% & 乾燥法 \\
            締固め度 & 90\%以上 & 密度試験 \\
            pH & 11.5以上 & pH試験 \\
            一軸圧縮強度（28日） & 400 kN/m$^2$以上 & 一軸圧縮試験 \\
            CBR & 30\%以上 & CBR試験 \\
            \bottomrule
          \end{tabular}
        \end{center}
      \end{block}
    \end{minipage}
  \end{center}
  
  \vspace{0.5em}
  
  \begin{columns}
    \begin{column}{0.48\linewidth}
      \begin{alertblock}{施工管理}
        \begin{itemize}
          \item 混合の均一性確保
          \item 養生期間の設定
          \item 環境条件の考慮
        \end{itemize}
      \end{alertblock}
    \end{column}
    \begin{column}{0.48\linewidth}
      \begin{exampleblock}{品質検査}
        \begin{itemize}
          \item 定期的なサンプリング
          \item 迅速試験法の活用
          \item データの蓄積と分析
        \end{itemize}
      \end{exampleblock}
    \end{column}
  \end{columns}
\end{frame}

\begin{frame}
  \frametitle{実用化への課題と対策}
  \begin{columns}
    \begin{column}{0.48\linewidth}
      \begin{alertblock}{技術的課題}
        \begin{enumerate}
          \item 長期耐久性の検証
          \item 多様な土質への適用性
          \item 施工の効率化
          \item コストの最適化
        \end{enumerate}
      \end{alertblock}
    \end{column}
    
    \begin{column}{0.48\linewidth}
      \begin{exampleblock}{対策と展望}
        \begin{enumerate}
          \item 促進劣化試験の実施
          \item 適用範囲の明確化
          \item 混合装置の改良
          \item スラグの安定供給体制
        \end{enumerate}
      \end{exampleblock}
    \end{column}
  \end{columns}
  
  \vspace{1em}
  
  \begin{center}
    \begin{minipage}{0.7\linewidth}
      \begin{block}{実用化に向けたロードマップ}
        \centering
        \small
        基礎研究 $\to$ 試験施工 $\to$ 実証試験 $\to$ 普及展開
      \end{block}
    \end{minipage}
  \end{center}
\end{frame}

%%%%%%%%%%%%%%%%%%%%%%%%%
% Environmental Safety
%%%%%%%%%%%%%%%%%%%%%%%%%
\section{環境安全性評価}

\begin{frame}
  \frametitle{環境影響評価}
  \begin{columns}
    \begin{column}{0.55\linewidth}
      \begin{block}{溶出試験結果}
        \begin{itemize}
          \item 重金属類：\emph{全て基準値以下}
          \item pH：環境基準を満足
          \item 有害物質：検出されず
        \end{itemize}
      \end{block}
      
      \vspace{1em}
      
      \begin{alertblock}{安全性の確認}
        \begin{itemize}
          \item 土壌汚染のリスクなし
          \item 地下水への影響なし
          \item 生態系への悪影響なし
        \end{itemize}
      \end{alertblock}
    \end{column}
    
    \begin{column}{0.42\linewidth}
      \begin{exampleblock}{環境貢献}
        \centering
        \Large
        \emph{循環型社会}\\
        \emph{の実現}
        
        \vspace{1em}
        \normalsize
        \begin{itemize}
          \item 産業副産物の有効利用
          \item CO$_2$排出量の削減
          \item 天然資源の保全
        \end{itemize}
      \end{exampleblock}
    \end{column}
  \end{columns}
  
  \vspace{0.5em}
  \centering
  \begin{minipage}{0.65\linewidth}
    \begin{roundbox}[green!25]{\linewidth}
      \centering
      \small
      環境と経済の\emph{両立}を実現する技術
    \end{roundbox}
  \end{minipage}
\end{frame}

%%%%%%%%%%%%%%%%%%%%%%%%%
% Conclusion
%%%%%%%%%%%%%%%%%%%%%%%%%
\section{結論}

\begin{frame}
  \frametitle{本研究の成果}
  \begin{center}
    \begin{minipage}{0.85\linewidth}
      \begin{block}{主要な成果}
        \begin{enumerate}
          \item 鉄鋼スラグによる土質改良効果を\emph{定量的}に解明
          \begin{itemize}
            \item 一軸圧縮強度：最大2.2倍向上
            \item CBR値：最大6倍向上
          \end{itemize}
          
          \item カルシア改質の\emph{メカニズム}を科学的に解明
          \begin{itemize}
            \item pH上昇と水和物生成の関係
            \item 微視的構造変化の観察
          \end{itemize}
          
          \item 実用化のための\emph{品質管理手法}を確立
          \begin{itemize}
            \item 管理基準値の設定
            \item 試験方法の標準化
          \end{itemize}
          
          \item \emph{環境安全性}を確認
          \begin{itemize}
            \item 溶出試験による安全性検証
            \item 環境負荷の評価
          \end{itemize}
        \end{enumerate}
      \end{block}
    \end{minipage}
  \end{center}
\end{frame}

\begin{frame}
  \frametitle{今後の展望}
  \begin{columns}
    \begin{column}{0.48\linewidth}
      \begin{alertblock}{学術的展開}
        \begin{itemize}
          \item より詳細な反応機構の解析
          \item 他の産業副産物への応用
          \item 新規改質材料の開発
          \item 理論モデルの構築
        \end{itemize}
      \end{alertblock}
    \end{column}
    
    \begin{column}{0.48\linewidth}
      \begin{exampleblock}{実用的展開}
        \begin{itemize}
          \item 実構造物への適用
          \item 施工技術の確立
          \item コスト削減手法の開発
          \item 標準仕様書への反映
        \end{itemize}
      \end{exampleblock}
    \end{column}
  \end{columns}
  
  \vspace{1.5em}
  
  \begin{center}
    \begin{minipage}{0.75\linewidth}
      \begin{block}{}
        \centering
        \Large
        \emph{持続可能な社会の実現}に向けた\\
        地盤工学技術の発展に貢献
      \end{block}
    \end{minipage}
  \end{center}
\end{frame}

%%%%%%%%%%%%%%%%%%%%%%%%%
% References
%%%%%%%%%%%%%%%%%%%%%%%%%
\begin{frame}[allowframebreaks]
  \frametitle{参考文献}
  \printbibliography[heading=none]
\end{frame}

%%%%%%%%%%%%%%%%%%%%%%%%%
% Acknowledgement
%%%%%%%%%%%%%%%%%%%%%%%%%
\begin{frame}
  \frametitle{謝辞}
  \begin{center}
    \begin{minipage}{0.85\linewidth}
      本研究を進めるにあたり、終始懇切なるご指導とご鞭撻を賜りました指導教員の○○教授に心より感謝申し上げます。
      
      \vspace{1em}
      
      また、貴重なご助言をいただきました副査の先生方、実験にご協力いただいた研究室の皆様、ならびに鉄鋼スラグを提供していただいた企業の皆様に深く感謝いたします。
      
      \vspace{2em}
      
      \Large
      \emph{ありがとうございました}
    \end{minipage}
  \end{center}
\end{frame}

% Appendix
\appendix

\begin{frame}
  \frametitle{付録：試験条件詳細}
  \begin{block}{供試体の作製条件}
    \begin{itemize}
      \item 供試体寸法：直径50 mm $\times$ 高さ100 mm
      \item 締固めエネルギー：E$_\mathrm{c}$ = 550 kJ/m$^3$
      \item 養生条件：温度20$\pm$2$^\circ$C、相対湿度95\%以上
      \item 繰返し回数：各条件につき5回以上
    \end{itemize}
  \end{block}
  
  \begin{block}{化学分析条件}
    \begin{itemize}
      \item XRD：CuK$\alpha$線、40 kV、30 mA
      \item SEM：加速電圧15 kV、金蒸着
      \item 溶出試験：環境庁告示46号に準拠
    \end{itemize}
  \end{block}
\end{frame}

\end{document}
